% ============================================================================
% CUDA 学习重点
% ============================================================================

\subsection{线程层级结构}

CUDA 的线程组织为两层结构：
\begin{itemize}
    \item \textbf{网格（Grid）}：由多个线程块（Block）组成
    \item \textbf{线程块（Block）}：由多个线程（Thread）组成
\end{itemize}

每个线程可以通过以下内置变量识别自己的位置：
\begin{itemize}
    \item \texttt{threadIdx.x} / \texttt{threadIdx.y} / \texttt{threadIdx.z}：线程在块内的索引
    \item \texttt{blockIdx.x} / \texttt{blockIdx.y} / \texttt{blockIdx.z}：块在网格内的索引
    \item \texttt{blockDim.x} / \texttt{blockDim.y} / \texttt{blockDim.z}：每个块的维度（大小）
    \item \texttt{gridDim.x} / \texttt{gridDim.y} / \texttt{gridDim.z}：网格的维度（大小）
\end{itemize}

\subsection{Hello CUDA 示例}

我们来看一个简单的 CUDA 程序：

\lstinputlisting[style=cstyle]{../CUDA/src/hello.cu}

代码说明：
\begin{itemize}
    \item \texttt{\_\_global\_\_} 表示这是一个在 GPU 上运行的内核函数
    \item \texttt{hello\_kernel<<<1, 4>>>()} 启动内核，参数表示：
        \begin{itemize}
            \item 网格大小：1 个块
            \item 块大小：4 个线程
        \end{itemize}
    \item \texttt{cudaDeviceSynchronize()} 同步主机和设备，等待所有 GPU 线程执行完毕
\end{itemize}

预期输出：
\begin{verbatim}
Hello from CPU
Hello from GPU, thread 0
Hello from GPU, thread 1
Hello from GPU, thread 2
Hello from GPU, thread 3
\end{verbatim}

\subsection{内存模型}

CUDA 设备有多层内存层次：
\begin{itemize}
    \item \textbf{全局内存（Global Memory）}：所有线程都可访问，容量大但速度较慢
    \item \textbf{共享内存（Shared Memory）}：块内线程共享，速度快但容量小
    \item \textbf{局部内存（Local Memory）}：每个线程私有
    \item \textbf{常量内存（Constant Memory）}：只读，所有线程共享
    \item \textbf{纹理内存（Texture Memory）}：只读，为特定访问模式优化
    \item \textbf{寄存器（Registers）}：每个线程私有，速度最快但数量有限
\end{itemize}

合理的内存使用是优化 CUDA 程序性能的关键。
